\chapter{Código del Proyecto}
\label{ch:codigo}

Este anexo presenta los fragmentos de código más relevantes del sistema de predicción de activos financieros desarrollado en Python.

\section{Módulo de Descarga de Datos}

El siguiente código muestra la clase encargada de descargar los datos de precios y sentimiento:

\lstset{style=codestyle, language=Python}
\begin{lstlisting}[frame=single, caption={Descargador de datos financieros}]
import os
import time
from datetime import datetime

import pandas as pd
import requests
import yfinance as yf

class DataDownloader:
    def __init__(self, data_dir="data/raw"):
        self.data_dir = data_dir
        os.makedirs(self.data_dir, exist_ok=True)
        
    def download_financial_data(self, tickers, start_date, end_date):
        """Descarga datos OHLCV de Yahoo Finance."""
        print(f"Downloading data for {tickers} from {start_date} to {end_date}...")
        for ticker in tickers:
            try:
                df = yf.download(
                    ticker,
                    start=start_date,
                    end=end_date,
                    auto_adjust=True,
                    progress=False,
                )
                if not df.empty:
                    filename = f"{ticker.replace('-', '_')}_history.csv"
                    filepath = os.path.join(self.data_dir, filename)
                    df.to_csv(filepath)
                    print(f"Saved {ticker} data to {filepath}")
                else:
                    print(f"No data found for {ticker}")
            except Exception as e:
                print(f"Error downloading {ticker}: {e}")
	    
    def _fetch_gdelt_timeline(self, query, mode, start_dt, end_dt, max_retries=3):
        url = "https://api.gdeltproject.org/api/v2/doc/doc"
        params = {
            "query": query,
            "mode": mode,
            "format": "json",
            "startdatetime": start_dt.strftime("%Y%m%d%H%M%S"),
            "enddatetime": end_dt.strftime("%Y%m%d%H%M%S"),
        }
        for attempt in range(1, max_retries + 1):
            try:
                response = requests.get(url, params=params, timeout=60)
                response.raise_for_status()
                try:
                    data = response.json()
                except ValueError:
                    preview = response.text[:200].strip().replace("\n", " ")
                    raise ValueError(f"GDELT non-JSON response: {preview}")
                timeline = data.get("timeline", [])
                if not timeline:
                    return pd.DataFrame()
                rows = timeline[0].get("data", [])
                if not rows:
                    return pd.DataFrame()
                df = pd.DataFrame(rows)
                df["date"] = pd.to_datetime(df["date"], format="%Y%m%dT%H%M%SZ", errors="coerce")
                df = df.dropna(subset=["date"]).set_index("date").sort_index()
                return df
            except Exception as e:
                if attempt >= max_retries:
                    raise e
                time.sleep(5 * attempt)
        return pd.DataFrame()

    def download_gdelt_sentiment(self, query="bitcoin", start_date=None, end_date=None,
                                 output_filename="sentiment_BTC_gdelt.csv"):
        """Descarga sentimiento desde GDELT (Average Tone + volumen)."""
        if not start_date or not end_date:
            print("start_date/end_date required for GDELT sentiment.")
            return False

        print(f"Downloading GDELT tone for '{query}'...")
        tone_frames = []
        volume_frames = []

        try:
            for chunk_start, chunk_end in self._iter_year_ranges(start_date, end_date):
                tone_df = self._fetch_gdelt_timeline(query, "timelinetone", chunk_start, chunk_end)
                if not tone_df.empty:
                    tone_df = tone_df.rename(columns={"value": "Sentiment_Index"})
                    tone_frames.append(tone_df)

                volume_df = self._fetch_gdelt_timeline(query, "timelinevol", chunk_start, chunk_end)
                if not volume_df.empty:
                    volume_df = volume_df.rename(columns={"value": "Sentiment_Volume"})
                    volume_frames.append(volume_df)

                time.sleep(2)

            if not tone_frames:
                print(f"No GDELT tone data for '{query}'.")
                return False

            tone = pd.concat(tone_frames)
            tone = tone[~tone.index.duplicated(keep="last")].sort_index()

            if volume_frames:
                volume = pd.concat(volume_frames)
                volume = volume[~volume.index.duplicated(keep="last")].sort_index()
                tone = tone.join(volume, how="left")

            filepath = os.path.join(self.data_dir, output_filename)
            tone.to_csv(filepath)
            print(f"Saved GDELT sentiment to {filepath}")
            return True
        except Exception as e:
            print(f"Error downloading GDELT sentiment: {e}")
            return False

    def _iter_year_ranges(self, start_date, end_date):
        start_dt = datetime.strptime(start_date, "%Y-%m-%d")
        end_dt = datetime.strptime(end_date, "%Y-%m-%d")
        year = start_dt.year
        while year <= end_dt.year:
            chunk_start = datetime(year, 1, 1)
            chunk_end = datetime(year, 12, 31)
            if year == start_dt.year:
                chunk_start = start_dt
            if year == end_dt.year:
                chunk_end = end_dt
            yield chunk_start, chunk_end
            year += 1
\end{lstlisting}

\section{Módulo de Preprocesamiento}

El preprocesador calcula indicadores técnicos y prepara los datos para el modelado:

\begin{lstlisting}[frame=single, caption={Preprocesador con indicadores técnicos}]
import numpy as np
import pandas as pd
from ta.trend import SMAIndicator, MACD
from ta.momentum import RSIIndicator
from ta.volatility import BollingerBands

class DataPreprocessor:
    def add_technical_indicators(self, df):
        """Anade indicadores tecnicos al DataFrame."""
        df['SMA_20'] = SMAIndicator(close=df['Close'], window=20).sma_indicator()
        df['SMA_50'] = SMAIndicator(close=df['Close'], window=50).sma_indicator()
        df['RSI'] = RSIIndicator(close=df['Close'], window=14).rsi()
        
        macd = MACD(close=df['Close'])
        df['MACD'] = macd.macd()
        df['MACD_Signal'] = macd.macd_signal()
        
        bb = BollingerBands(close=df['Close'], window=20, window_dev=2)
        df['BB_High'] = bb.bollinger_hband()
        df['BB_Low'] = bb.bollinger_lband()
        return df

    def merge_sentiment(self, price_df, sentiment_df):
        """Une precios y sentimiento sin look-ahead (forward fill con valores pasados)."""
        merged = price_df.join(sentiment_df, how='left')
        return merged.ffill()

    def create_target(self, df, horizon=1):
        """Target: rendimiento logarítmico y dirección a h días."""
        df[f'Target_Return_{horizon}d'] = np.log(df['Close'].shift(-horizon) / df['Close'])
        df[f'Target_Direction_{horizon}d'] = (df[f'Target_Return_{horizon}d'] > 0).astype(int)
        return df
\end{lstlisting}

\section{Modelo GRU}

Arquitectura del modelo GRU, una de las redes recurrentes evaluadas:

\begin{lstlisting}[frame=single, caption={Arquitectura del modelo GRU}]
from tensorflow.keras.models import Sequential
from tensorflow.keras.layers import GRU, Dense, Dropout, Input

def build_gru_model(input_shape):
    """Construye el modelo GRU para prediccion."""
    model = Sequential()
    model.add(Input(shape=input_shape))
    model.add(GRU(50, return_sequences=True))
    model.add(Dropout(0.2))
    model.add(GRU(50, return_sequences=False))
    model.add(Dropout(0.2))
    model.add(Dense(25))
    model.add(Dense(1))
    model.compile(optimizer='adam', loss='mean_squared_error')
    return model
\end{lstlisting}

\section{Normalización para Deep Learning}

Fragmento simplificado del escalado para evitar fuga de información. Se ajusta
el \texttt{MinMaxScaler} solo con entrenamiento y se aplica al test.

\begin{lstlisting}[frame=single, caption={Escalado sin fuga de información}]
from sklearn.preprocessing import MinMaxScaler

scaler = MinMaxScaler(feature_range=(0, 1))
target_scaler = MinMaxScaler(feature_range=(0, 1))

# Eliminar columnas con informacion futura (Target_*)
leaky_cols = [c for c in train_df.columns if str(c).startswith("Target_")]
train_X = train_df.drop(columns=leaky_cols)
test_X = test_df[train_X.columns]

# Ajuste solo en entrenamiento
scaled_train = scaler.fit_transform(train_X)
target_scaler.fit(train_X[['Close']])

# Aplicación en test (sin refit)
scaled_test = scaler.transform(test_X)
\end{lstlisting}

\section{Semillas de reproducibilidad}

Inicialización de semillas para repetir resultados de entrenamiento y evaluación.

\begin{lstlisting}[frame=single, caption={Semillas y determinismo}]
from src.utils.reproducibility import set_global_seed

SEED = 42
set_global_seed(SEED)
\end{lstlisting}

\section{Validación Walk-Forward}

Implementación del esquema de validación temporal:

\textit{Nota:} en los experimentos se instancia con \texttt{n\_splits=12}; aquí se muestra la clase base.

\begin{lstlisting}[frame=single, caption={Validador walk-forward}]
import numpy as np
from sklearn.metrics import mean_squared_error, mean_absolute_error
from scipy import stats

class TimeSeriesValidator:
    def __init__(self, n_splits=5, test_size=30):
        self.n_splits = n_splits
        self.test_size = test_size
    
    def split(self, df):
        """Genera splits temporales para validacion."""
        n_samples = len(df)
        if n_samples < self.test_size * (self.n_splits + 1):
            raise ValueError("No hay datos suficientes.")
        fold_size = n_samples // (self.n_splits + 1)
        
        for i in range(1, self.n_splits + 1):
            train_end = fold_size * i
            test_end = min(train_end + self.test_size, n_samples)
            if test_end <= train_end:
                break
            
            train_df = df.iloc[:train_end]
            test_df = df.iloc[train_end:test_end]
            
            yield train_df, test_df
    
    def evaluate(self, y_true, y_pred, y_prev=None):
        """Calcula metricas de evaluacion."""
        rmse = np.sqrt(mean_squared_error(y_true, y_pred))
        mae = mean_absolute_error(y_true, y_pred)
        mask = y_true != 0
        mape = np.mean(np.abs((y_true[mask] - y_pred[mask]) / y_true[mask])) * 100
        
        if y_prev is not None and len(y_prev) == len(y_true):
            actual_change = y_true - y_prev
            pred_change = y_pred - y_prev
        else:
            actual_change = np.diff(y_true) if len(y_true) > 1 else np.array([])
            pred_change = np.diff(y_pred) if len(y_pred) > 1 else np.array([])
        
        if len(actual_change) > 0:
            mask = np.isfinite(actual_change) & np.isfinite(pred_change)
            actual_change = actual_change[mask]
            pred_change = pred_change[mask]
            direction_mask = actual_change != 0
            actual_dir = (actual_change[direction_mask] > 0).astype(int)
            pred_dir = (pred_change[direction_mask] > 0).astype(int)
        else:
            actual_dir = np.array([])
            pred_dir = np.array([])
        
        if len(actual_dir) > 0:
            dir_acc = np.mean(actual_dir == pred_dir) * 100
            tp = int(np.sum((actual_dir == 1) & (pred_dir == 1)))
            tn = int(np.sum((actual_dir == 0) & (pred_dir == 0)))
            fp = int(np.sum((actual_dir == 0) & (pred_dir == 1)))
            fn = int(np.sum((actual_dir == 1) & (pred_dir == 0)))
            denom_f1 = (2 * tp + fp + fn)
            dir_f1 = (2 * tp / denom_f1) * 100 if denom_f1 > 0 else 0.0
            denom_mcc = np.sqrt((tp + fp) * (tp + fn) * (tn + fp) * (tn + fn))
            dir_mcc = (tp * tn - fp * fn) / denom_mcc if denom_mcc > 0 else 0.0
        else:
            dir_acc = np.nan
            dir_f1 = np.nan
            dir_mcc = np.nan
        return {'RMSE': rmse, 'MAE': mae, 'MAPE': mape,
                'DirectionalAccuracy': dir_acc, 'DirectionalF1': dir_f1, 'DirectionalMCC': dir_mcc}
\end{lstlisting}

\section{Script Principal de Ejecución}

El pipeline completo se orquesta desde el script principal:

\begin{lstlisting}[frame=single, caption={Script de ejecucion principal}]
import os
import sys
import shutil
from pathlib import Path


def sync_plots_to_memoria(src_dir="results/plots", dst_dir="memoria/Imagenes"):
    src_path = Path(src_dir)
    dst_path = Path(dst_dir)
    if not src_path.exists():
        return

    dst_path.mkdir(parents=True, exist_ok=True)

    for plot_path in src_path.glob("*.png"):
        shutil.copy2(plot_path, dst_path / plot_path.name)

def main():
    project_root = os.path.dirname(os.path.abspath(__file__))
    existing_pythonpath = os.environ.get("PYTHONPATH", "")
    if existing_pythonpath:
        os.environ["PYTHONPATH"] = project_root + os.pathsep + existing_pythonpath
    else:
        os.environ["PYTHONPATH"] = project_root

    os.environ.setdefault("PYTHONDONTWRITEBYTECODE", "1")
    print("TFG: Predicción Financiera - Ejecución Completa", flush=True)

    # 1. Descarga de Datos
    print("\n[1/6] Descargando datos...", flush=True)
    if os.system(f"{sys.executable} src/data/downloader.py") != 0:
        raise SystemExit("Error en la descarga de datos.")
    
    # 2. Preprocesamiento
    print("\n[2/6] Procesando datos...", flush=True)
    if os.system(f"{sys.executable} src/data/preprocessor.py") != 0:
        raise SystemExit("Error en el preprocesamiento.")
    
    # 3. Ejecutar Experimentos Principales
    print("\n[3/6] Ejecutando experimentos (A-D)...", flush=True)
    if os.system(f"{sys.executable} scripts/run_all_experiments.py") != 0:
        raise SystemExit("Error en los experimentos.")
    
    # 4. Métricas adicionales
    print("\n[4/6] Calculando métricas direccionales y test Diebold-Mariano...", flush=True)
    if os.system(f"{sys.executable} scripts/run_additional_metrics.py") != 0:
        raise SystemExit("Error calculando metricas adicionales.")

    # 5. Generar gráficas
    print("\n[5/6] Generando gráficas (y sincronizando para LaTeX)...", flush=True)
    if os.system(f"{sys.executable} scripts/generate_all_plots.py") != 0:
        raise SystemExit("Error generando graficas.")
    if os.system(f"{sys.executable} scripts/generate_summary_plots.py") != 0:
        raise SystemExit("Error generando graficas de resumen.")
    sync_plots_to_memoria()
    
    # 6. Exportar tablas LaTeX
    print("\n[6/6] Exportando tablas a LaTeX...", flush=True)
    if os.system(f"{sys.executable} scripts/export_tables_to_latex.py") != 0:
        raise SystemExit("Error exportando tablas LaTeX.")

    print("\nEjecución Finalizada", flush=True)
    print("Revisa results/ para los CSV y results/plots/ para las figuras.", flush=True)
    print("Para compilar la memoria, las figuras también quedan en memoria/Imagenes/.", flush=True)
    
if __name__ == "__main__":
    main()
\end{lstlisting}

\section{Métricas Adicionales de Evaluación}

Método para calcular el test de Diebold-Mariano:

\begin{lstlisting}[frame=single, caption={Test de Diebold-Mariano}]
from scipy import stats
import numpy as np

class TimeSeriesValidator:
    def diebold_mariano_test(self, errors_1, errors_2, horizon=1):
        """
        Test de Diebold-Mariano para comparar precisión predictiva de dos modelos.
        H0: ambos modelos tienen igual precisión.
        """
        d = errors_1**2 - errors_2**2
        d = d[np.isfinite(d)]

        n = len(d)
        if n < 2:
            return {"dm_statistic": np.nan, "p_value": np.nan,
                    "significant": False, "better_model": None}

        mean_d = np.mean(d)

        # Varianza de largo plazo con corrección Newey-West
        gamma_0 = np.sum((d - mean_d)**2) / n
        gamma_sum = 0
        for k in range(1, min(horizon, n)):
            gamma_k = np.sum((d[:-k] - mean_d) * (d[k:] - mean_d)) / n
            gamma_sum += gamma_k

        var_mean_d = (gamma_0 + 2 * gamma_sum) / n
        if var_mean_d <= 0 or not np.isfinite(var_mean_d):
            var_mean_d = np.var(d, ddof=1) / n
        if var_mean_d <= 0:
            return {"dm_statistic": np.nan, "p_value": np.nan,
                    "significant": False, "better_model": None}

        dm_stat = mean_d / np.sqrt(var_mean_d)
        p_value = 2 * (1 - stats.t.cdf(abs(dm_stat), df=n-1))

        return {
            "dm_statistic": dm_stat,
            "p_value": p_value,
            "significant": p_value < 0.05,
            "better_model": 1 if mean_d < 0 else 2,
        }
\end{lstlisting}

En los scripts del trabajo, el test se invoca con \texttt{horizon=7} para mantener coherencia con el horizonte principal de predicción.

El test de Wilcoxon para la ablación se implementa en:

\noindent\texttt{scripts/run\_all\_experiments.py}

Los resultados se guardan en:

\noindent\texttt{results/experiment\_c\_ablation\_tests.csv}.

El código fuente completo del proyecto está disponible en el repositorio del trabajo.
